\section*{THESIS SUMMARY}
\thispagestyle{plain}

This bachelor thesis studies Orthogonal Time Frequency Space (OTFS) modulation and its extension with Index Modulation (OTFS-IM) for high-mobility wireless communication scenarios. The main motivation is that conventional Orthogonal Frequency Division Multiplexing (OFDM) is sensitive to Doppler-induced channel variation, while OTFS represents information in the delay-Doppler domain, where multipath delay and Doppler effects can be described more naturally.

The thesis first develops a baseline OTFS transceiver and then extends it to OTFS-IM. In OTFS-IM, information is conveyed through both QAM constellation symbols and the indices of active positions inside each delay-Doppler index-modulation block. This creates a design trade-off among bit error rate, spectral efficiency, active-position selection, and detector reliability. To study this trade-off, the thesis implements Orthogonal Hamming Distance (OHD) and Balanced OHD (B-OHD) pattern selection and compares them with baseline OTFS, OFDM-LMMSE, and a random-pattern diagnostic reference.

MATLAB simulations are conducted under a discrete delay-Doppler multipath Rayleigh fading channel with additive white Gaussian noise. The evaluation focuses on BER, index BER, symbol BER, pattern error rate, and BER-spectral-efficiency trade-off. The results show that OTFS is more robust than OFDM-LMMSE in the considered high-Doppler channel model. OTFS-IM can further improve BER in selected configurations, but its benefit must be interpreted together with spectral efficiency and the chosen \((n,k)\) activation structure.

\textbf{Chapter 1 -- Introduction}

This chapter presents the motivation, application context, research problem, objectives, scope, and methodology of the thesis. It explains the difficulty of high-mobility wireless communication, the limitations of OFDM under Doppler effects, and the reason for studying OTFS and Index Modulation as a combined transmission approach.

\textbf{Chapter 2 -- Theoretical Background}

This chapter reviews the theoretical foundation of the thesis, including wireless fading channels, multipath propagation, Doppler effects, doubly selective channels, OFDM, delay-Doppler representation, and Index Modulation. These concepts provide the basis for the OTFS and OTFS-IM system models developed in the following chapters.

\textbf{Chapter 3 -- OTFS System Design and Modeling}

This chapter develops the baseline OTFS transceiver. It describes the transmitter, delay-Doppler grid, ISFFT, Heisenberg transform, discrete delay-Doppler channel, receiver processing chain, equivalent input-output model, and message passing detector. This baseline is used as the reference system before index modulation is introduced.

\textbf{Chapter 4 -- OTFS with Index Modulation Design}

This chapter extends the baseline OTFS system to OTFS-IM. It presents block-wise index modulation, index-bit and symbol-bit mapping, active-position pattern tables, receiver-side MAP decision, and the OHD and B-OHD pattern-selection methods. The chapter also discusses random-pattern validation and the complexity of exact pattern search.

\textbf{Chapter 5 -- Simulation Results and Performance Evaluation}

This chapter evaluates OFDM-LMMSE, baseline OTFS, OHD OTFS-IM, and B-OHD OTFS-IM under common channel and noise assumptions. The analysis includes BER curves, error decomposition, pattern-selection validation, configuration-wise BER comparison, and BER-SE trade-off across multiple \((n,k)\) configurations.

\textbf{Chapter 6 -- Conclusion and Future Work}

This chapter summarizes the main contributions and findings of the thesis. It concludes that OTFS is a strong candidate for high-Doppler delay-Doppler channels and that OTFS-IM provides a flexible BER-SE design space through active-position selection. Remaining limitations and future extensions, including velocity-based channel modeling, coded OTFS-IM, adaptive pattern selection, and improved detectors, are also discussed.

\cleardoublepage
